\documentclass[11pt,a4paper]{article}

%--------------------------------------------------
% PACKAGES
%--------------------------------------------------
\usepackage[utf8]{inputenc}
\usepackage[T1]{fontenc}
\usepackage[french]{babel}
\addto\captionsfrench{%
\renewcommand{\contentsname}{\color{esieadarkblue}\sffamily Sommaire}%
}
\usepackage{lmodern}
\usepackage{geometry}
\usepackage{graphicx}
\usepackage{hyperref}
\usepackage{setspace}
\usepackage{parskip}
\usepackage{xcolor}
\usepackage{fancyhdr}
\usepackage{tikz}
\usepackage{titlesec}
\usepackage{xurl}
\usepackage{enumitem}
\usepackage{tcolorbox}
\usepackage{mdframed}
\usepackage{fontawesome5}
\usepackage{booktabs}
\usepackage{array}
\usepackage{microtype}
\usepackage{tocloft}

% Sections et sous-sections dans le sommaire
\setcounter{tocdepth}{2}

% Style des entrées
\renewcommand{\cftsecfont}{\sffamily\bfseries\color{headerblue}}
\renewcommand{\cftsecpagefont}{\sffamily\bfseries\color{esieadarkblue}}
\renewcommand{\cftsecleader}{\cftdotfill{\cftdotsep}}
\setlength{\cftbeforesecskip}{6pt}
\setlength{\cftsecindent}{0em}
\setlength{\cftsecnumwidth}{2em}

\tcbuselibrary{skins, breakable, most}

%--------------------------------------------------
% COULEURS
%--------------------------------------------------
\definecolor{esieablue}{RGB}{54,169,225}
\definecolor{esieadarkblue}{RGB}{21,29,52}
\definecolor{esieaorange}{RGB}{245,158,0}
\definecolor{headerblue}{RGB}{0,68,158}
\definecolor{pourquoibg}{RGB}{255,255,255}
\definecolor{commentbg}{RGB}{255, 255, 255}
\definecolor{pourquoiborder}{RGB}{54,169,225}
\definecolor{commentborder}{RGB}{245,158,0}
\definecolor{sectionrule}{RGB}{0,68,158}
\definecolor{lightgray}{RGB}{243,248,250}
\definecolor{warningbg}{RGB}{255,248,230}
\definecolor{warningborder}{RGB}{245,158,0}

%--------------------------------------------------
% MISE EN PAGE
%--------------------------------------------------
\geometry{
	margin=2.5cm,
	top=3cm,
	bottom=3.5cm,
	footskip=2cm
}
\setlength{\parskip}{0.6em}
\setlength{\parindent}{0pt}

%--------------------------------------------------
% EN-TÊTE / PIED DE PAGE
%--------------------------------------------------
\pagestyle{fancy}
\fancyhf{}
\setlength{\headheight}{22pt}

\fancyhead[L]{%
	\colorbox{esieadarkblue}{%
		\makebox[\headwidth][c]{%
			\vspace{3pt}%
			\textcolor{white}{\small\textbf{\textsf{RESTITUTION DIAGNOSTIC MONAIDECYBER}}}%
			\vspace{3pt}%
		}%
	}%
}

\fancyfoot[L]{\small\textcolor{esieadarkblue}{\textbf{MonAideCyber ID~<<DIAGNOSTIC_ID>>}}}
\fancyfoot[R]{\small\thepage}
\renewcommand{\headrulewidth}{0pt}
\renewcommand{\footrulewidth}{0.4pt}

%--------------------------------------------------
% TITRES DE SECTIONS
%--------------------------------------------------
\titleformat{\section}
{\color{headerblue}\normalfont\LARGE\bfseries\sffamily}
{\thesection}{1em}{}
[\vspace{-0.3em}\textcolor{headerblue}{\rule{\linewidth}{1.2pt}}]

\titleformat{\subsection}
{\color{esieadarkblue}\normalfont\Large\bfseries\sffamily}
{\thesubsection}{1em}{}
[\vspace{-0.5em}\textcolor{esieablue}{\rule{0.4\linewidth}{0.6pt}}]

\titlespacing*{\section}{0pt}{1.5em}{0.8em}
\titlespacing*{\subsection}{0pt}{1.2em}{0.5em}

%--------------------------------------------------
% BOÎTES COLORÉES
%--------------------------------------------------
\newtcolorbox{pourquoibox}{
	breakable,
	enhanced,
	colback=pourquoibg,
	colframe=pourquoiborder,
	leftrule=4pt, toprule=0.5pt, bottomrule=0.5pt, rightrule=0.5pt,
	arc=2pt,
	boxsep=4pt,
	left=6pt, right=6pt, top=4pt, bottom=4pt,
	title={\faQuestionCircle\enskip\textbf{Pourquoi ?}},
	fonttitle=\color{pourquoiborder}\sffamily\bfseries,
	coltitle=pourquoiborder,
	attach boxed title to top left={yshift=-2mm, xshift=6pt},
	boxed title style={colback=pourquoibg, colframe=pourquoiborder, arc=2pt, leftrule=4pt}
}

\newtcolorbox{commentbox}{
	breakable,
	enhanced,
	colback=commentbg,
	colframe=commentborder,
	leftrule=4pt, toprule=0.5pt, bottomrule=0.5pt, rightrule=0.5pt,
	arc=2pt,
	boxsep=4pt,
	left=6pt, right=6pt, top=4pt, bottom=4pt,
	title={\faTools\enskip\textbf{Comment mettre en \oe{}uvre cette mesure ?}},
	fonttitle=\color{commentborder}\sffamily\bfseries,
	coltitle=commentborder,
	attach boxed title to top left={yshift=-2mm, xshift=6pt},
	boxed title style={colback=commentbg, colframe=commentborder, arc=2pt, leftrule=4pt}
}

%--------------------------------------------------
% LISTES
%--------------------------------------------------
\setlist[itemize]{
	label=\textcolor{esieablue}{\textbullet},
	leftmargin=1.5em,
	itemsep=0.3em,
	topsep=0.3em
}
\setlist[enumerate]{
	leftmargin=1.5em,
	itemsep=0.3em,
	topsep=0.3em
}

%--------------------------------------------------
% HYPERLIENS
%--------------------------------------------------
\hypersetup{
	colorlinks=true,
	linkcolor=esieadarkblue,
	urlcolor=headerblue,
	pdfauthor={ESIEA},
	pdftitle={Restitution Diagnostic MonAideCyber}
}

%--------------------------------------------------
% COMMANDES UTILES
%--------------------------------------------------
\newcommand{\recommandation}[3]{%
	\subsection{#1}
	\begin{pourquoibox}
		#2
	\end{pourquoibox}
	\begin{commentbox}
		#3
	\end{commentbox}
	\vspace{0.5em}
}

\newcommand{\sectionruleline}{%
	\vspace{0.3em}
	\textcolor{sectionrule}{\rule{\linewidth}{0.8pt}}
	\vspace{0.3em}
}

%==================================================
\begin{document}
	%==================================================
	
	%--------------------------------------------------
	% TABLE DES MATIÈRES
	%--------------------------------------------------
	{
		\hypersetup{linkcolor=esieadarkblue}
		{\setlength{\parskip}{0pt}\tableofcontents}
	}
	\thispagestyle{fancy}
	\clearpage
	
	
	%--------------------------------------------------
	% REMERCIEMENTS
	%--------------------------------------------------
	\section*{Remerciements}
	\addcontentsline{toc}{section}{Remerciements}
	
	\begin{tcolorbox}[
		colback=lightgray,
		colframe=esieadarkblue,
		arc=4pt,
		boxrule=0.8pt,
		left=10pt, right=10pt, top=8pt, bottom=8pt
		]
		Nous tenons à vous remercier sincèrement pour l'accueil que vous nous avez réservé lors de notre récente visite dans le cadre du diagnostic de premier niveau au sein de votre société.
		
		Votre disponibilité, la qualité des échanges ainsi que la transparence dont vous avez fait preuve ont grandement facilité le bon déroulement de cette mission. Grâce à votre collaboration, nous avons pu recueillir des informations précieuses et mieux comprendre vos activités, vos enjeux ainsi que vos attentes.
		
		Nous espérons que ce premier diagnostic pourra constituer une base utile pour la suite de vos démarches, et restons bien entendu à votre disposition pour tout complément ou accompagnement futur.
		
		Encore merci pour votre confiance et votre engagement dans cette démarche.
	\end{tcolorbox}
	
	\clearpage
	
	%--------------------------------------------------
	% INDICATEURS
	%--------------------------------------------------
	\section*{Indicateurs}
	\addcontentsline{toc}{section}{Indicateurs MonAideCyber}
	
	<<INDICATEUR_POLAIRE>>
	
	\clearpage
	
	%--------------------------------------------------
	% I) RECOMMANDATIONS PRIORITAIRES
	%--------------------------------------------------
	\section{Recommandations prioritaires}
	\label{sec:recommandations-prioritaires}
	
	<<MESURES_PRIORITAIRES>>
	
	\newpage
	
	
	%--------------------------------------------------
	% II) AUTRES RECOMMANDATIONS — NON TECHNIQUE
	%--------------------------------------------------
	<<SECTION_MESURES_COMPLEMENTAIRES>>
	
	\clearpage
	
	%--------------------------------------------------
	% III) LIENS UTILES
	%--------------------------------------------------
	\section{Liens utiles}
	\label{sec:liens-utiles}
	
	\begin{tcolorbox}[
		breakable,
		enhanced,
		colback=lightgray,
		colframe=esieadarkblue,
		arc=4pt,
		boxrule=0.8pt,
		title={\sffamily\bfseries\color{white} Contacts et ressources utiles},
		colbacktitle=esieadarkblue,
		fonttitle=\sffamily\bfseries,
		left=6pt, right=6pt, top=4pt, bottom=4pt
		]
		\small
		
		\textbf{\faLock\enskip MonExpertCyber}\\
		Service de mise en relation avec des prestataires labellisés ExpertCyber.\\
		\href{https://cybermalveillance.gouv.fr/accompagnement}{https://cybermalveillance.gouv.fr/accompagnement}
		
		\vspace{3pt}\sectionruleline
		
		\textbf{\faPhone\enskip 17Cyber}\\
		Service public gratuit d'assistance en ligne destiné aux victimes de cybermalveillance.\\
		\href{https://cybermalveillance.gouv.fr/17cyber}{https://cybermalveillance.gouv.fr/17cyber}
		
		\vspace{3pt}\sectionruleline
		
		\textbf{\faServer\enskip CSIRT — Centres de réponse à incidents}\\
		Les CSIRT régionaux sont des centres de réponse à incidents cyber au profit des entités du territoire régional.\\
		\href{https://cert.ssi.gouv.fr/csirt/csirt-regionaux}{https://cert.ssi.gouv.fr/csirt/csirt-regionaux}
		
		\vspace{3pt}\sectionruleline
		
		\textbf{\faUniversity\enskip Ma sécurité, service de l'État}\\
		La gendarmerie et la police nationale vous accompagnent dans vos démarches de sécurité cyber.\\
		\href{https://masecurite.interieur.gouv.fr/fr}{https://masecurite.interieur.gouv.fr/fr}
		
		\vspace{3pt}\sectionruleline
		
		\textbf{\faBuilding\enskip L'ANSSI}\\
		L'Agence nationale de la sécurité des systèmes d'information.\\
		\href{https://cyber.gouv.fr}{https://cyber.gouv.fr}\\
		Délégués régionaux de l'ANSSI : \href{https://cyber.gouv.fr/action-territoriale}{https://cyber.gouv.fr/action-territoriale}\\
		Produits et services évalués : \href{https://cyber.gouv.fr/trouver-un-produitservice-de-securite-evalue}{https://cyber.gouv.fr/trouver-un-produitservice-de-securite-evalue}
		
		\vspace{3pt}\sectionruleline
		
		\textbf{\faBook\enskip Ressources MonAideCyber}\\
		\textit{Les CGU} :\\
		\href{https://monaide.cyber.gouv.fr/cgu}{https://monaide.cyber.gouv.fr/cgu}\\[0.2em]
		\textit{L'équipe MonAideCyber} — Remarques, questions ou remontées après votre diagnostic :\\
		\href{mailto:MonaideCyber@et.esiea.fr}{MonaideCyber@et.esiea.fr}
		\par\noindent\textcolor{esieadarkblue}{\rule{\linewidth}{0.8pt}}\par

	\end{tcolorbox}
	
\end{document}
